% Section 4 -- Results. Numbers are macros and \input{tables/...} blocks that
% `python train.py paper` fills from results/.
\section{Results}
\label{sec:results}

\subsection{Setup}
All experiments use the cleaned dataset and the fixed group-aware 5-fold split of
Section~\ref{sec:method_data}. Training ran with PyTorch and timm under
automatic mixed precision. Folds 0--3 were trained on an NVIDIA T4 (fp16) and
fold 4 on an NVIDIA RTX 4090 (bf16). Data split, hyper-parameters, schedule and
seed are identical; only data-loader workers and activation checkpointing,
which affect speed but not the computation, were changed. Hyper-parameters are listed in
Section~\ref{sec:method_twostage} and in \texttt{configs/asmoe\_v2.yaml} of the
accompanying repository. Every number is the mean $\pm$ SD over five folds.
Tables were generated automatically on \ReportGenerated.

\subsection{Main comparison}
\label{sec:results_main}

\input{tables/main_results}

Table~\ref{tab:main_results} compares ASMOE~V2 with the two strongest legacy
models, re-evaluated under the same leakage-free protocol.
ASMOE~V2 reaches \VtwoAcc\,$\pm$\,\VtwoAccSD\,\% accuracy and
\VtwoFOne\,$\pm$\,\VtwoFOneSD\,\% Macro-F1 (MCC \VtwoMCC), a change of
\DeltaFOne\,pp Macro-F1 and \DeltaAcc\,pp accuracy over the legacy MobileViT-S
baseline (\BaseFOne\,$\pm$\,\BaseFOneSD\,\% Macro-F1). Pooled over the five
validation folds, the number of errors falls from 464 (of 3{,}076 images) to
222 (of 3{,}012 cleaned images). ASMOE~V2 does not reach the
$\geq95\,\%$ accuracy and Macro-F1 target set at the start of this work. The
remaining gap is concentrated in Phytophthora and in the Fungi/Pest pair
(Section~\ref{sec:results_perclass}). The improvement is larger than the
fold-to-fold variation: the weakest ASMOE~V2 fold (89.05\,\% Macro-F1, 91.36\,\%
accuracy) is above the strongest legacy fold (85.15\,\% and 86.34\,\%), and the
SD of accuracy falls from 1.27 to 0.91\,pp. Because data cleaning,
input resolution, training recipe and architecture all changed together, the gain
reflects the complete pipeline rather than the mixture alone
(Section~\ref{sec:discussion}).

\begin{figure}[h]
  \centering
  \safegraphic[width=\textwidth]{figures/learning_curves.png}
  \caption{ASMOE~V2 validation Macro-F1 and loss per epoch (mean $\pm$ SD over
  five folds). The dashed line marks the end of the 5-epoch frozen-backbone
  stage. Compare with the legacy curves in Fig.~\ref{fig:legacy_curves}.}
  \label{fig:v2_curves}
\end{figure}

Figure~\ref{fig:v2_curves} shows that training converges: validation loss
flattens after about 30 of the 45 epochs, and Macro-F1 changes by less than
1\,pp over the last 15 epochs, so a longer schedule is unlikely to close the
gap to the target.

\subsection{Per-class recovery of minority classes}
\label{sec:results_perclass}

\input{tables/per_class}

\begin{figure}[h]
  \centering
  \safegraphic{figures/per_class_f1.png}
  \caption{Per-class F1 of the legacy baseline and ASMOE V2 (mean over 5 folds).}
  \label{fig:per_class}
\end{figure}

The legacy model's weakest classes were Nematode (\BaseFOneNematode\,\% F1),
Pest (\BaseFOnePest\,\%) and Healthy (\BaseFOneHealthy\,\%). ASMOE~V2 reaches
\VtwoFOneNematode\,\% (\DeltaFOneNematode\,pp) on Nematode,
\VtwoFOnePest\,\% (\DeltaFOnePest\,pp) on Pest and
\VtwoFOneHealthy\,\% (\DeltaFOneHealthy\,pp) on Healthy
(Table~\ref{tab:per_class}, Fig.~\ref{fig:per_class}).
In the legacy model, Nematode and Healthy errors were mainly false positives
caused by loss re-weighting (Section~\ref{sec:audit_imbalance}). The gain on
these two classes comes mainly from higher precision. On the pooled matrices,
Nematode precision rises from 67.0 to 82.3\,\% (false positives $30\rightarrow14$)
and recall from 89.7 to 97.0\,\%. Healthy precision rises from 75.2 to 86.8\,\%
(false positives $57\rightarrow29$) and recall from 86.1 to 94.5\,\%. Precision
remains the limiting factor for both classes, so over-prediction of minority
classes is reduced but not removed. The supervised contrastive term and the
milder re-weighting were introduced together, so their individual contributions
cannot be separated here. Removing the six label-conflicting duplicates, all of which
involved Pest, removes errors that no model could avoid, so part of the Pest
gain is a measurement correction rather than a modelling improvement.
Re-training the legacy MobileViT-S on the cleaned split with its original recipe
would separate the two effects.

\begin{figure}[h]
  \centering
  \safegraphic[width=0.8\textwidth]{figures/asmoe_v2_cm.png}
  \caption{Pooled 5-fold confusion matrix of ASMOE V2. Compare with
  Fig.~\ref{fig:legacy_cm}.}
  \label{fig:v2_cm}
\end{figure}

Figure~\ref{fig:v2_cm} should be read against the legacy matrix
(Fig.~\ref{fig:legacy_cm}). Pest$\rightarrow$Fungi errors fall from 76 to 27 and
Fungi$\rightarrow$Pest errors from 69 to 19. Virus$\rightarrow$Healthy errors
fall from 20 to 8. Phytophthora$\rightarrow$Fungi is the only major confusion
that does not improve (31 to 38). Consistent with this, Phytophthora has the
smallest F1 gain (\DeltaFOnePhytopthora\,pp) and is now the weakest class
(\VtwoFOnePhytopthora\,\%). Closing the Phytophthora/Fungi gap is the most direct
route to higher Macro-F1.

\subsection{Routing behaviour}
\label{sec:results_routing}

\begin{figure}[h]
  \centering
  \safegraphic{figures/router_weights.png}
  \caption{Mean routing weight per expert, grouped by true class (validation folds).}
  \label{fig:router}
\end{figure}

Figure~\ref{fig:router} shows the average router output for each true class. If
the experts specialise as intended, the global-attention expert should receive
more weight for Virus and Healthy, where the evidence is diffuse, and the
texture expert more weight for Pest and Fungi. The observed weights only
partly follow this hypothesis. Swin-B receives its largest mean weights on Virus
and Pest (0.45 each), but DenseNet-201 is used mainly for Bacteria (0.51) and
receives only 0.17 on Pest and 0.27 on Fungi. EfficientNetV2-M dominates Healthy
(0.50) and Nematode (0.52). The routing is therefore clearly class-dependent:
the DenseNet-201 weight ranges from 0.09 (Nematode) to 0.51 (Bacteria), and the
fold-to-fold SD of every class-mean weight is at most 0.07. However, the learned
division of labour is not the texture-versus-global split we expected. The
routing is therefore input-dependent, but this study does not measure how much
accuracy it adds over a fixed average of the same experts.

\subsection{Explainability}
\label{sec:results_xai}

\begin{figure}[p]
  \centering
  \safegraphic[width=0.85\textwidth]{figures/gradcam_00_Bacteria.png}\\[2pt]
  \safegraphic[width=0.85\textwidth]{figures/gradcam_01_Fungi.png}\\[2pt]
  \safegraphic[width=0.85\textwidth]{figures/gradcam_02_Healthy.png}\\[2pt]
  \safegraphic[width=0.85\textwidth]{figures/gradcam_03_Nematode.png}\\[2pt]
  \safegraphic[width=0.85\textwidth]{figures/gradcam_04_Pest.png}\\[2pt]
  \safegraphic[width=0.85\textwidth]{figures/gradcam_05_Phytopthora.png}\\[2pt]
  \safegraphic[width=0.85\textwidth]{figures/gradcam_06_Virus.png}
  \caption{Routing-weighted Grad-CAM for one correctly classified validation image
  per class (fold 0). Each row shows the input, the fused map with the predicted
  class and confidence, and the maps of the local-texture (DenseNet-201),
  multi-scale (EfficientNetV2-M) and global-attention (Swin-B) experts with their
  routing weights $w_e$ for that image.}
  \label{fig:gradcam}
\end{figure}

Routing-weighted Grad-CAM (Fig.~\ref{fig:gradcam}) shows which pixels each expert
uses and how much each expert contributes. The expert maps differ in character.
The EfficientNetV2-M maps are compact and centred on individual lesions or leaf
regions. The DenseNet-201 maps cover larger leaf areas when the expert receives a
high weight (Bacteria, $w=0.67$) and are almost flat when its weight is low. The
Swin-B maps are patchy and, in several examples (Pest, Virus), also highlight
background regions at the image border. Because the fused map is the
weight-averaged sum, these background responses carry through when Swin-B
dominates the routing. The lesion-focus ratio (fraction of CAM mass inside the leaf mask,
computed by \texttt{train.py gradcam}) quantifies this. On the eight fold-0
examples, a mean of 79\,\% of the fused CAM mass lies inside the leaf mask. Six
of the eight examples exceed 0.68, and three exceed 0.93. The exception is a Pest
image with a ratio of 0.31, where most attention falls outside the colour-based
leaf mask. Eight images are too few for a general statement, but they indicate
that the decision usually rests on foliage rather than soil or background.
